\chapter*{常见缩写和术语先导}
\phantomsection
\addcontentsline{toc}{chapter}{常见缩写和术语先导}

本卷从第一章开始就会同时出现 C++、ELF、ABI、section、symbol、PMU、benchmark、GTest 等词。不要把它们当成孤立名词背下来；每个词都对应代码中的一个对象、命令行输出中的一个字段，或报告中的一条证据。

\begin{description}
  \item[CPU: Central Processing Unit] 中央处理器。它执行机器指令、访问寄存器和内存，并通过缓存、流水线、分支预测和乱序执行影响程序速度。
  \item[ELF: Executable and Linkable Format] Linux 常见的可执行文件和目标文件格式。本卷只实现 ELF64 小端切片，重点训练 header、section、string table 和 symbol table 的证据链。
  \item[ABI: Application Binary Interface] 应用二进制接口。它约定函数参数、返回值、寄存器保存、栈对齐、符号和链接边界，让不同编译单元能在二进制层面协作。
  \item[ISA: Instruction Set Architecture] 指令集架构。x86-64、AArch64、RISC-V 都是 ISA；ELF header 的 machine 字段能告诉我们二进制面向哪类机器。
  \item[x86-64] 64 位 x86 指令集。本卷第一版以 x86-64 为主要观察对象，但解析器会把 AArch64 和 RISC-V machine 字段也分类出来。
  \item[section] ELF 文件中的一段逻辑区域，例如 \code{.text}、\code{.rodata}、\code{.symtab} 和 \code{.strtab}。section 是链接器和调试工具常用的组织方式。
  \item[symbol] 符号。函数名、全局对象或链接器可见标记都可能成为符号。符号把源码名字、section、地址、大小和绑定关系连接起来。
  \item[string table] 字符串表。ELF 中很多名字不直接写在结构体里，而是保存一个偏移，偏移指向某个字符串表。
  \item[entry address] 入口地址。程序开始执行的虚拟地址，不等于 C++ 的 \code{main} 函数地址；运行时启动代码会先执行。
  \item[CLI: Command-Line Interface] 命令行接口。本卷的 \code{cpu1ee\_cli} 用文件路径和模式参数输出 report 或 section CSV。
  \item[CSV: Comma-Separated Values] 逗号分隔表格。适合保存 section 摘要、benchmark 结果和可脚本复查的字段。
  \item[checksum] 校验摘要。本卷使用它固定输入字节身份，避免报告只描述文件名而没有描述内容。
  \item[reference] 参考实现。可以慢，但必须确定、清楚、可复查。实践中先有 reference，再让优化或工具扩展与它对齐。
  \item[oracle] 判定器。它比较 reference 和待测实现，判断输出是否满足同一语义。
  \item[contract] 合同。代码边界的明确约定：输入是什么，输出是什么，错误怎样表达，哪些行为暂不支持。
  \item[GTest: GoogleTest] C++ 单元测试框架。本卷用它固定 ELF fixture、错误输入、符号上限和格式化输出。
  \item[GMock: GoogleMock] GoogleTest 配套 mock 框架。复杂工程可用它替换外部依赖；本卷先通过 GTest/GMock 工程接入建立测试能力。
  \item[Google Benchmark] C++ benchmark 库。本卷用它做解析路径 smoke，不用 benchmark 替代正确性测试。
  \item[PMU: Performance Monitoring Unit] 性能监控单元。它统计 retired instructions、cache miss、branch miss 等硬件事件，是性能报告的重要证据来源。
  \item[perf] Linux 性能分析工具，可读取 PMU 事件、采样热点和统计执行行为。没有 perf 权限时，报告要写明限制。
  \item[objdump] 二进制工具，可反汇编、查看符号和 section 信息。它常用来把 C++ 函数和机器指令对应起来。
  \item[readelf] ELF 检查工具，可查看 ELF header、section table、symbol table 等字段。它是本卷解析器的交叉验证对象。
  \item[Debug/Release] 常见构建类型。Debug 偏调试信息和可读性；Release 打开优化。两者的机器码和性能结论不能混用。
\end{description}
